\chapter{EMDA Code Setup on MaryLou}

\section{Building \texttt{emda-gr}}
To build \texttt{emda-gr}, the following external packages are required:
\begin{itemize}
\item \textbf{C/C++11 or higher} (tested with GNU and Intel compilers)
\item \textbf{MPI implementation} (tested with \texttt{MVAPICH2}, \texttt{OpenMPI}, and Intel MPI)
\item \textbf{BLAS/LAPACK} (tested with \texttt{OpenBLAS} and Intel MKL)
\item \textbf{GSL} (GNU Scientific Library); set \texttt{GSL\_ROOT\_DIR} in \texttt{CMake} to enable automatic detection
\end{itemize}

On MaryLou, the recommended module environment is:

\begin{verbatim}
andrewc0@login03:$ ml
Currently Loaded Modules:

1. gcc/14.1
2. python/3.12
3. cmake/3.28
4. gcc-runtime/13.2.0-g4hc7ev
5. eigen/3.4.0-6kqviac
6. openmpi5-GR
7. gsl-GR
8. mkl-GR
9. rclone/1.65.2-tyg4lvy
   \end{verbatim}

To ensure that these dependencies are always available, it is convenient to load
them automatically in your shell configuration. Open your
\texttt{\textasciitilde/.bash\_profile}:

\begin{verbatim}
vim ~/.bash_profile
\end{verbatim}

An example \texttt{.bash\_profile} is shown below. Additional customizations may
be added as needed.

\begin{verbatim}

# --- Modules ---

module purge
module use /grphome/grp_GR/.modulefiles
module load gcc/14.1 python/3.12 cmake/3.28 eigen/3.4.0-6kqviac
module load openmpi5-GR gsl-GR mkl-GR rclone

export HCOLL_DIR=/vapps/rhel9/x86_64/HPC-X/hpcx-v2.19-gcc-inbox-redhat9-cuda12-x86_64/hcoll/lib
export LD_LIBRARY_PATH="${HCOLL_DIR}:${LD_LIBRARY_PATH:-}"

# --- GSL ---

export GSL_DIR=/grphome/grp_GR/apps/gsl-2.8
export GSL_ROOT_DIR=/grphome/grp_GR/apps/gsl-2.8
export LD_LIBRARY_PATH="${GSL_DIR}/lib:${LD_LIBRARY_PATH:-}"

# --- MKL ---

if [ -n "${MKLROOT:-}" ] && [ -d "${MKLROOT}/lib/intel64" ]; then
export LD_LIBRARY_PATH="${MKLROOT}/lib/intel64:${LD_LIBRARY_PATH:-}"
fi

# --- GCC 14 runtime: libstdc++ and libgcc_s ---

# Required to avoid GLIBCXX and CXXABI runtime errors.

if command -v g++ >/dev/null 2>&1; then
GCC_LIBDIR="$(dirname "$(g++ -print-file-name=libstdc++.so.6)")"
GCC_GCCLIBDIR="$(dirname "$(g++ -print-file-name=libgcc_s.so.1)")"

if [ -d "$GCC_LIBDIR" ]; then
export LD_LIBRARY_PATH="$GCC_LIBDIR:${LD_LIBRARY_PATH:-}"
fi
if [ -d "$GCC_GCCLIBDIR" ]; then
export LD_LIBRARY_PATH="$GCC_GCCLIBDIR:${LD_LIBRARY_PATH:-}"
fi
fi

# --- History settings ---

HISTSIZE=1000000
HISTFILESIZE=2000000
HISTCONTROL=ignoredups:erasedups
HISTIGNORE="ls:bg:fg:history"
shopt -s histappend

# Save history after every command.

PROMPT_COMMAND="history -a; history -c; history -r; $PROMPT_COMMAND"
\end{verbatim}

After reloading the shell, verify that the modules have been loaded correctly:

\begin{verbatim}
ml
\end{verbatim}

The output should match the module list shown above.

\section{Building the Solver}

First, clone the EMDA repository. Create a project directory if desired, then run:

\begin{verbatim}
git clone [git@github.com](mailto:git@github.com):dfvankomen/emda-gr.git
cd emda-gr
\end{verbatim}

CMake will automatically fetch \texttt{dendrolib} during configuration.

Create and enter a build directory:

\begin{verbatim}
mkdir build
cd build
\end{verbatim}

Configure the project using:

\begin{verbatim}
ccmake ..
\end{verbatim}

CMake will locate the required libraries and open an interactive configuration
interface. The main settings that may need to be adjusted are the system of
equations and the target CPU architecture.

The available equation systems include EMDA, EMD, and EMGR. For a first test
run, choose the EMDA system.

The CPU architecture should match the compute nodes on which the job will run.
On MaryLou, the \texttt{physics2} and \texttt{m12} partitions use 64-core AMD
EPYC 7763 processors. These are Zen~3 CPUs, so set

\begin{verbatim}
CPU_ARCH = znver3
\end{verbatim}

After configuration is complete, generate the build files from within
\texttt{ccmake}. Then compile the code using:

\begin{verbatim}
make -j
\end{verbatim}

\subsection{Build Outputs}

Once the build completes, the \texttt{emdaSolver} and \texttt{tpid} binaries
should be located in the solver build directory, for example:

\begin{verbatim}
emda-gr/build/emdaSolver/
\end{verbatim}

At this point, the solver is ready to be used for simulations.

\section{Runs and Job Scripts}

After building the solver, move from the home directory to the scratch directory.
The location of the scratch directory depends on the supercomputer. On MaryLou,
the scratch area is located under \texttt{/nobackup/autodelete/usr/}.

First, change into your scratch area:

\begin{verbatim}
cd /nobackup/autodelete/usr/YOURUSERNAME/
\end{verbatim}

Create a directory for the EMDA runs and enter it:

\begin{verbatim}
mkdir EMDA_runs
cd EMDA_runs
\end{verbatim}

Next, copy the solver binaries and a parameter file into the run directory. In
the commands below, replace \texttt{\textasciitilde/path/to/} with the location
of your local \texttt{emda-gr} checkout.

Copy \texttt{emdaSolver}:

\begin{verbatim}
cp ~/path/to/emda-gr/build/emdaSolver/emdaSolver ./
\end{verbatim}

Copy \texttt{tpid}:

\begin{verbatim}
cp ~/path/to/emda-gr/build/emdaSolver/tpid ./
\end{verbatim}

Copy the parameter file:

\begin{verbatim}
cp ~/path/to/emda-gr/emda_simplified.param.toml ./
\end{verbatim}

Verify that everything copied correctly:

\begin{verbatim}
ls
\end{verbatim}

You should see:

\begin{verbatim}
emda_simplified.param.toml  emdaSolver  tpid
\end{verbatim}

Now create the standard output directories for the run:

\begin{verbatim}
mkdir cp prof vtu bah
\end{verbatim}

These directories store the main simulation outputs:
\begin{itemize}
\item \texttt{cp}: checkpoint files, which allow the run to be restarted
\item \texttt{prof}: waveform extraction, ADM quantities, charges, and constraint outputs
\item \texttt{vtu}: visualization files that can be opened in ParaView
\item \texttt{bah}: apparent-horizon data
\end{itemize}

A later section will describe the most important parameter-file settings needed
to control these outputs.


\subsection{Submitting a Run with Slurm}

To run the simulation on the cluster, create a Slurm job script. For example:

\begin{verbatim}
vim Job.sb
\end{verbatim}

Paste the following into the file:

\begin{verbatim}
#!/bin/bash
#SBATCH --time=72:00:00
#SBATCH --ntasks=256
#SBATCH --nodes=2
#SBATCH --mem-per-cpu=4096M
#SBATCH -J EMDA_RUN
#SBATCH --mail-user=[youremail@email.com](mailto:youremail@email.com)
#SBATCH --mail-type=END,FAIL
#SBATCH --partition=m12

echo "Starting run"

source ~/.bash_profile

./tpid emda_simplified.param.toml 256 &> TPID.txt
mpirun -np 256 ./emdaSolver emda_simplified.param.toml &>> Inspiral.txt

echo "Run completed"
\end{verbatim}

The \texttt{tpid} command generates the initial data, while
\texttt{emdaSolver} evolves that data forward in time. Since these commands are
written on separate lines and the \texttt{tpid} command is not placed in the
background, the script waits for \texttt{tpid} to finish before starting
\texttt{emdaSolver}. This ordering is necessary because \texttt{tpid} produces
the binary initial-data file, \texttt{rit\_q2\_tpid\_sol.bin}, which is then
read by \texttt{emdaSolver}.

On the \texttt{m12} partition, each node has 128 CPU cores. Therefore, a run
using 256 MPI tasks typically uses 2 nodes. If the number of MPI tasks is
changed, the number of requested nodes should be adjusted accordingly.

Submit the job with:

\begin{verbatim}
sbatch Job.sb
\end{verbatim}

Once submitted, the job will enter the queue and begin running when resources
become available. You can check the status of your jobs with:

\begin{verbatim}
squeue | grep yourusername
\end{verbatim}

After the job starts, you can monitor the output by inspecting the text files
generated during execution. For example:

\begin{verbatim}
cat TPID.txt
\end{verbatim}

or

\begin{verbatim}
tail -f Inspiral.txt
\end{verbatim}

At this point, the simulation is running.

\section{Parameter-File Explanation}

The parameter file is read by both \texttt{tpid} and \texttt{emdaSolver}. It
allows many run settings to be changed without rebuilding the code, which makes
it easy to adjust output frequencies, grid settings, checkpointing behavior,
and initial-data parameters.

This section summarizes some of the most important parameters. Parameters not
discussed here can usually be left at their default values.

\begin{verbatim}

# Whether or not the solver should attempt to restore from a checkpoint

EMDA_RESTORE_SOLVER = false
\end{verbatim}

This option determines whether \texttt{emdaSolver} should attempt to restart
from an existing checkpoint. For a fresh run, this should be set to
\texttt{false}. To continue a previous run from checkpoint data, set this to
\texttt{true}.

\begin{verbatim}
EMDA_IO_OUTPUT_FREQ = 400
\end{verbatim}

This controls how often VTU output is written. Smaller values produce output
more frequently, but they also increase storage usage. This parameter should be
chosen based on how much visualization data is needed.

\begin{verbatim}
EMDA_REMESH_TEST_FREQ = 200
\end{verbatim}

This sets how often the solver invokes the adaptive mesh refinement algorithm
to determine whether the grid should be refined or coarsened. If this value is
too small, remeshing can occur too frequently and may introduce unnecessary grid
noise. If it is too large, the grid may not adapt quickly enough to resolve the
dynamics. A value around \texttt{200} has worked well in typical EMDA runs.

\begin{verbatim}
EMDA_GW_EXTRACT_FREQ = 25
\end{verbatim}

This controls how often gravitational-wave quantities are extracted. The same
output step also includes other useful quantities, such as black-hole positions
and masses. This output is not especially expensive compared with VTU output,
so values in the range \texttt{25}--\texttt{100} are usually reasonable.

\begin{verbatim}
EMDA_CHECKPT_FREQ = 100
\end{verbatim}

This specifies the number of time steps between checkpoints. Checkpointing is
relatively inexpensive compared with full visualization output, so this value
can be kept fairly small. More frequent checkpoints make it easier to restart a
run after a crash or wall-time limit.

\begin{verbatim}

# VTU file prefix

EMDA_VTU_FILE_PREFIX = "vtu/emda_gr"

# Checkpoint file prefix

EMDA_CHKPT_FILE_PREFIX = "cp/emda_cp"

# Profiling file prefix

EMDA_PROFILE_FILE_PREFIX = "prof/emda_prof"
\end{verbatim}

These options control where different types of output data are written. The
\texttt{vtu} directory stores visualization output, the \texttt{cp} directory
stores checkpoint files, and the \texttt{prof} directory stores extracted
waveforms and other reduced data products.

\begin{verbatim}

# Which evolution variables to output

EMDA_VTU_OUTPUT_EVOL_INDICES = [0, 1, 2, 3, 4, 5, 6, 7, 8, 9,
10, 11, 12, 13, 14, 15, 16, 17, 18, 19, 20, 21, 22, 23, 24,
25, 26, 27, 28, 29, 30, 31, 32, 33, 34, 35]

# Number of evolution variables to output

EMDA_NUM_EVOL_VARS_VTU_OUTPUT = 36
\end{verbatim}

These settings determine which evolution variables are written to the VTU
files. These variables include quantities such as the metric components,
extrinsic-curvature components, gauge variables, and matter fields. If storage
usage becomes a concern, some variables may be omitted from the VTU output.

\begin{verbatim}

# Which constraint variables to output

EMDA_VTU_OUTPUT_CONST_INDICES = [0, 1, 2, 3, 4, 5, 6, 7,
8, 9, 10, 11, 12, 13]

# Number of constraint variables to output

EMDA_NUM_CONST_VARS_VTU_OUTPUT = 14
\end{verbatim}

These options are analogous to the evolution-variable output settings, but they
apply to constraint quantities. These include quantities such as the Hamiltonian
constraint, the momentum constraints, and additional constraints associated
with the electromagnetic sector.

\begin{verbatim}

# Whether to isolate an X slice

EMDA_VTU_X_SLICE = false

# Whether to isolate a Y slice

EMDA_VTU_Y_SLICE = false

# Whether to isolate a Z slice

EMDA_VTU_Z_SLICE = true

# Any combination may be set to true

\end{verbatim}

Writing VTU output for the full 3-D grid can be expensive, so it is often
preferable to output a 2-D slice instead. For most visualization purposes, a
single slice is sufficient. In typical runs, the Z slice is a good default
choice.

Although multiple slice options can be enabled at the same time, doing so may
produce output that is less straightforward to interpret. For simple
visualization, enable only one slice direction at a time.

\begin{verbatim}
EMDA_DENDRO_GRAIN_SZ = 50
\end{verbatim}

This controls the approximate number of grid blocks assigned to each processor.
In practice, this parameter should usually be kept between \texttt{25} and
\texttt{200}. Very small or very large values may reduce parallel efficiency.

\begin{verbatim}
EMDA_DENDRO_AMR_FAC = 0.2
\end{verbatim}

This controls the sensitivity of the adaptive mesh refinement algorithm. Smaller
values generally make the grid more sensitive to refinement, while larger
values make refinement less aggressive. This parameter should be adjusted with
care, since it affects both accuracy and computational cost.

\begin{verbatim}
EMDA_ELE_ORDER = 6
\end{verbatim}

This sets the finite-element order used by the Dendro infrastructure. Loosely
speaking, it controls how much information each block must exchange with its
neighbors. Higher values increase communication costs and may slow the code. For
the current EMDA implementation, this should generally be left at
\texttt{6}.

\begin{verbatim}

# Minimum depth of the mesh/octree

EMDA_MINDEPTH = 4

# Maximum depth of the mesh/octree

EMDA_MAXDEPTH = 15
\end{verbatim}

These parameters control the minimum and maximum refinement levels of the
octree mesh. The minimum depth sets the coarsest allowed grid resolution, while
the maximum depth sets the finest allowed grid resolution. Increasing
\texttt{EMDA\_MAXDEPTH} improves the maximum available resolution, but it also
increases memory usage and computational cost.

The CFL condition is set by the finest level of the grid, so increasing
\texttt{EMDA\_MAXDEPTH} can significantly increase the runtime. Each increase
in \texttt{EMDA\_MAXDEPTH} halves the grid spacing on the finest level. As a
rule of thumb, each black hole should be resolved by at least roughly 50 points
across the horizon. For an equal-mass binary with total mass $M=1$ on a
$\pm 400$ grid, \texttt{EMDA\_MAXDEPTH = 15} is usually adequate. For larger
mass ratios, the smaller black hole requires additional refinement. For example,
a $q=2$ binary typically requires \texttt{EMDA\_MAXDEPTH = 16}.

\begin{verbatim}
########################################################################

# REFINEMENT STRATEGY

# Which refinement mode to use

# 0 - WAMR

# 1 - EH

# 2 - EH_WAMR

# 3 - BH_LOC

# 4 - BH_WAMR

EMDA_REFINEMENT_MODE = 4
\end{verbatim}

This parameter controls the refinement strategy used by the code. The
recommended choice is currently \texttt{EMDA\_REFINEMENT\_MODE = 4}, which uses
black-hole-centered refinement together with wavelet adaptive mesh refinement.
\begin{verbatim}
EMDA_AMR_R_RATIO = 1.618033988749895
\end{verbatim}

This parameter determines how quickly the forced refinement around each black
hole falls off with radius. The default value is reasonable and should usually
be left unchanged for initial runs.

\begin{verbatim}

# Wavelet Tolerance Function to use (see grUtils.cpp for adding more)

# 1,2 - deprecated

# 3 - center highly refined; low refinement elsewhere eases over time

# 4 - start with refinement centered at BHs, then widen to r=0 focused

# 5 - focus refinement in self-similar ways about each BH

# 6 - refine well in r <= t and coarsely in r > t

EMDA_USE_WAVELET_TOL_FUNCTION = 6
\end{verbatim}

This parameter controls how the wavelet refinement tolerance varies across the
computational domain. Function~6, developed by William, has worked well in
practice and is the recommended choice for current EMDA runs. More details can
be found in Ref.~\cite{Black2025_ORBIT}.

\begin{verbatim}

# Minimum wavelet tolerance

EMDA_WAVELET_TOL = 0.00001

# Gravitational wave refinement tolerance

EMDA_GW_REFINE_WTOL = 0.0001

# Maximum wavelet tolerance

EMDA_WAVELET_TOL_MAX = 0.001

# Wavelet tolerance radii (for shell-based tolerance)

EMDA_WAVELET_TOL_FUNCTION_R0 = 3.0
EMDA_WAVELET_TOL_FUNCTION_R1 = 12.0
\end{verbatim}

The minimum wavelet tolerance, \texttt{EMDA\_WAVELET\_TOL}, is the most important
of these parameters because it sets the strictest refinement criterion used by
the grid. Smaller values force more refinement and generally improve accuracy,
but they also increase memory usage and runtime. The remaining parameters
depend on the chosen wavelet tolerance function. The default values are
reasonable starting points, but they may be adjusted for higher-resolution runs
or for runs in which the wave-extraction region requires additional refinement.

\begin{verbatim}

# Indices of variables used for refinement (see grDef.h for definitions)

EMDA_REFINE_VARIABLE_INDICES = [0, 1, 2, 3, 4, 5, 6, 7, 8, 9,
10, 11, 12, 13, 14, 15, 16, 17, 18, 19, 20, 21, 22, 23, 24,
25, 26, 27, 28, 29, 30, 31, 32, 33, 34, 35]

# Number of variables used for refinement

EMDA_NUM_REFINE_VARS = 36
\end{verbatim}

These parameters specify which evolution variables are used when deciding
whether the grid should be refined. It is not always necessary to use every
evolution variable as a refinement variable, but the full list is a safe default
for general-purpose EMDA runs.

\begin{verbatim}

# CFL (Courant--Friedrichs--Lewy) factor

EMDA_CFL_FACTOR = 0.25

# Simulation start time

EMDA_RK_TIME_BEGIN = 0

# Simulation end time

EMDA_RK_TIME_END = 1000
\end{verbatim}

These parameters control the CFL factor and the total coordinate time of the
simulation. The CFL factor sets the time step relative to the grid spacing on
the finest refinement level. The parameters \texttt{EMDA\_RK\_TIME\_BEGIN} and
\texttt{EMDA\_RK\_TIME\_END} determine the start and end times of the
evolution. The remaining Runge--Kutta settings can usually be left at their
default values.

\begin{verbatim}
###########################################

# EMDA THEORY PARAMETERS

###########################################
EMDA_LAMBDA = [1, 1, 1, 1]
EMDA_ALPHA_THEORY = [1.0, 1.0]
EMDA_LF = [1.0, 0.0]
EMDA_P_EXPO = -1.0
EMDA_ALPHA_POS_NEG_FLOOR = false
EMDA_ETA = [0.4, 0.4]
EMDA_XI = [0, 0, 0]
CHI_FLOOR = 0.00001
EMDA_TRK0 = 0.0
EMDA_ETA_R0 = 1.31
EMDA_ETA_POWER = [2.0, 2.0]
EMDA_ETADAMP = 1.0

ETA_CONST = 2.0
ETA_R0 = 30.0
ETA_DAMPING = 1.0
ETA_DAMPING_EXP = 1.0
\end{verbatim}

These parameters define the theory and gauge choices used in the evolution. The
main parameters that are typically changed are the two theory couplings
$\alpha_0$ and $\alpha_1$, stored in \texttt{EMDA\_ALPHA\_THEORY}. The choices

\[
(\alpha_0,\alpha_1) = (1,1)
\]

correspond to EMDA,

\[
(\alpha_0,\alpha_1) = (1,0)
\]

correspond to EMD, and

\[
(\alpha_0,\alpha_1) = (0,0)
\]

correspond to Einstein--Maxwell. The remaining gauge parameters are generally
safe to leave at their default values.

\begin{verbatim}

# Kreiss--Oliger dissipation parameter

KO_DISS_SIGMA = 0.4
\end{verbatim}

This parameter controls the strength of Kreiss--Oliger numerical dissipation.
Some dissipation is needed for stability, and \texttt{KO\_DISS\_SIGMA = 0.4} is
a reasonable starting value.

\begin{verbatim}
########################################################################

# Select initial data type

# See the bottom of the file for reference

EMDA_ID_TYPE = 0
\end{verbatim}

This parameter selects the type of initial data used by the evolution. The
choice \texttt{EMDA\_ID\_TYPE = 0} uses the TPID solver. Other black-hole
initial-data options are available, but TPID is the standard choice for the
binary black-hole simulations described here.

\begin{verbatim}
########################################################################

# Set up EMDA grid extents

EMDA_GRID_MIN_X = -400.0
EMDA_GRID_MAX_X = 400.0
EMDA_GRID_MIN_Y = -400.0
EMDA_GRID_MAX_Y = 400.0
EMDA_GRID_MIN_Z = -400.0
EMDA_GRID_MAX_Z = 400.0
\end{verbatim}

These parameters define the physical extent of the computational domain. The
values shown here are reasonable defaults for many binary simulations. The outer
boundary should be placed far enough from the wave-extraction region that
boundary reflections do not contaminate the extracted waveforms during the time
interval of interest.

\begin{verbatim}
########################################################################

# Black hole parameters

EMDA_BH1_AMR_R = 1.0
EMDA_BH1_CONSTRAINT_R = 2.0
EMDA_BH1_MAX_LEV = 15
EMDA_BH1 = {MASS = 0.5, X = 5.0, Y = 0.0, Z = 0.0,
V_X = -0.00100282, V_Y = 0.095775311, V_Z = 0.0,
SPIN = 0.0, SPIN_THETA = 0.0, SPIN_PHI = 0.0,
ELECTRIC_CHARGE = 0.05, MAGNETIC_CHARGE = 0.0}

EMDA_BH2_AMR_R = 1.0
EMDA_BH2_CONSTRAINT_R = 2.0
EMDA_BH2_MAX_LEV = 15
EMDA_BH2 = {MASS = 0.5, X = -5.0, Y = 0.0, Z = 0.0,
V_X = 0.00100282, V_Y = -0.095775311, V_Z = 0.0,
SPIN = 0.0, SPIN_THETA = 0.0, SPIN_PHI = 0.0,
ELECTRIC_CHARGE = 0.05, MAGNETIC_CHARGE = 0.0}
\end{verbatim}

These parameters define the black-hole properties used by the evolution and by
the refinement infrastructure. The parameter \texttt{EMDA\_BH1\_AMR\_R} sets
the radius within which refinement is forced around the first black hole, while
\texttt{EMDA\_BH1\_CONSTRAINT\_R} sets the radius excluded when computing
constraint norms. The corresponding \texttt{BH2} parameters play the same role
for the second black hole. The parameters \texttt{EMDA\_BH1\_MAX\_LEV} and
\texttt{EMDA\_BH2\_MAX\_LEV} set the maximum refinement level allowed near each
black hole.

The \texttt{X}, \texttt{Y}, and \texttt{Z} entries do not directly determine the
black holes' initial positions in the TPID solve. Instead, they specify where
the evolution code begins searching for the black-hole locations and where it
centers the black-hole refinement regions. For equal-mass binaries, these
values should match the corresponding $\pm b$ positions used in the TPID
parameters.

The electric charge should remain below the black-hole mass in order to avoid
exceeding the extremal limit. Magnetic charge is still under development and
should not be used for standard production runs. The spin parameter corresponds
to the dimensionless spin $\chi$, with angular momentum $S=\chi m^2$.

TPID reads the charge and spin from these black-hole parameter blocks, but it
does not use the listed mass or position values in the same way as the
evolution code. The TPID mass and separation are set separately by the TPID
parameters described below.

Finally, note that \texttt{V\_X}, \texttt{V\_Y}, and \texttt{V\_Z} represent
the black-hole momenta, not their coordinate velocities. This distinction is
important when setting up binary initial data.

\begin{verbatim}
########################################################################

# GRAVITATIONAL WAVE EXTRACTION PARAMETERS

EMDA_GW_NUM_RADAII = 6
EMDA_GW_RADAII = [50.0, 60.0, 70.0, 80.0, 90.0, 100.0]

EMDA_GW_NUM_LMODES = 7
EMDA_GW_L_MODES = [2, 3, 4, 5, 6, 7, 8]

# Radius at which conserved quantities are extracted

EMDA_GW_EXTRACTION_RADIUS = 300.0

# Modes for EM, dilaton, and axion radiation

EMDA_RAD_NUM_LMODES = 6
EMDA_RAD_L_MODES = [1, 2, 3, 4, 5, 6]
########################################################################
\end{verbatim}

These parameters control the extraction of gravitational, electromagnetic,
dilaton, and axion radiation. Radiation is extracted on spherical surfaces at
the radii listed in \texttt{EMDA\_GW\_RADAII}. The gravitational-wave modes are
set by \texttt{EMDA\_GW\_L\_MODES}, while the modes for the electromagnetic,
dilaton, and axion fields are set by \texttt{EMDA\_RAD\_L\_MODES}.

The mode numbers correspond to the spherical-harmonic multipoles being
extracted. Although modes can be extracted up to roughly $\ell \approx 10$, the
dominant contributions are usually contained in the lower modes, especially
$\ell \leq 4$.

\begin{verbatim}
# Time steps between apparent-horizon (AEH) solves
AEH_SOLVER_FREQ = 50
\end{verbatim}

This parameter controls how often the apparent-horizon solver is called. The
remaining AEH-related settings should generally be left at their default values
unless the apparent-horizon finder fails or additional horizon information is
needed.

\begin{verbatim}
TPID_PAR_B = 5.0
TPID_TARGET_M_PLUS  = 0.48594
TPID_TARGET_M_MINUS = 0.48594
\end{verbatim}

The parameter \texttt{TPID\_PAR\_B} sets the coordinate half-separation used by
TPID. For example, \texttt{TPID\_PAR\_B = 5.0} corresponds to punctures located
near $x=\pm 5$. The target masses are the values used internally by TPID when
constructing the initial data.

\begin{verbatim}
TPID_CENTER_OFFSET = {X = 0.0, Y = 0.0, Z = 0.0}
\end{verbatim}

For equal-mass binaries, this should usually be left at zero. For unequal-mass
binaries, the system may need to be recentered so that the center of mass is
located near the origin.

\begin{verbatim}
TPID_NPOINTS_A   = 40
TPID_NPOINTS_B   = 40
TPID_NPOINTS_PHI = 20
\end{verbatim}

These parameters control the resolution of the TPID grid. The values shown here
are reasonable defaults for initial testing, but higher values may be needed for
more accurate initial data.

\begin{verbatim}
TPID_GIVE_BARE_MASS = 1
\end{verbatim}

This parameter determines whether the puncture masses are fixed or solved for.
When \texttt{TPID\_GIVE\_BARE\_MASS = 0}, TPID solves for the bare masses. When
\texttt{TPID\_GIVE\_BARE\_MASS = 1}, TPID uses the specified target masses.

A useful procedure is to first estimate the desired masses, run TPID with mass
solving enabled, then update the target masses and rerun with
\texttt{TPID\_GIVE\_BARE\_MASS = 1} to finalize the initial data.

\begin{verbatim}
TPID_REPLACE_LAPSE_WITH_SQRT_CHI = true
\end{verbatim}

This option should be enabled when using the slow-start lapse.

\begin{verbatim}
# Include dilatonic contributions in rho during TPID solve
TPID_USE_DILATON = true

# Interpolated initial data should not be used with TPID_USE_DILATON
# You should also be using the EMD equations in this case
TPID_USE_INTERPOLATED_SOLUTION = false
\end{verbatim}

These parameters control how scalar-field contributions are included in the
initial-data solve. If both options are set to \texttt{false}, TPID generates
Einstein--Maxwell initial data. Setting
\texttt{TPID\_USE\_INTERPOLATED\_SOLUTION = true} produces EMD-type initial
data. Setting \texttt{TPID\_USE\_DILATON = true} includes dilatonic
contributions directly in the TPID solve.

The initial-data choice should be consistent with both the theory parameters in
\texttt{EMDA\_ALPHA\_THEORY} and the system of equations selected when the code
was configured with CMake.
